% Article VII. Recovery by dependency, 10-09-2026.
% RESULTADO_RECUPERADO: active integral 2249, fuente/colaboracion/
% partes_iv_v/tex/residencias/86_cosmologia_cerrada.tex (S04),
% lines 2--44 (selected charts), 174--248 (tensor and acceleration).
% Source S04 read in full, SHA-256:
% feb2582c3316c8b30829cfb0476bfa2f4e95889dabfba16fe43f2d835a226673.
% Definition, uniqueness, conservation, decomposition, exchange, and
% acceleration are recovered with their explicit conditions.
% FORMALIZACION_REUNIDA: relation to the effective source in 31;
% total conservation is distinguished from separate matter conservation.
% CERTIFICADO_NUEVO: expanded proofs of trace, balance, and signs;
% local Ricci calculation for the three charts already selected by S04.
% This does not modify the bounce in 50 or determine an amplitude s_0 here.
%
% INTEGRATION DEPENDENCIES:
% - vii:sec:realizacion-cartan: the metric comes from the realized co-tetrad.
% - vii:sec:acoplamientos: G_int and the clock t_0 of the same HMT state;
%   c_int=1 is a choice of units subsequent to their construction.
% - vii:eq:ecuacion-metrica-efectiva and vii:eq:balance-metrico-total, from 31:
%   used only to identify the residual on a solution.
% - T_b is the tensor published by the declared matter transducer;
%   its separate conservation is an expressly stated condition.
% - The TRIT chart selector and t_0(x)>0 are antecedents of the de Sitter
%   specialization, not observational cosmological data.
% This fragment does not certify closure of VII as a whole.

\section{Residual tensor, trace decomposition, and covariant balance}
\label{vii:sec:tensor-residual-balance}

The geometric realization of the enriched state provides a metric,
while the matter transducer expresses the source to which the readout
refers. The residual tensor retains the difference between both
representations in a common tensor space. Its decomposition distinguishes
the sector proportional to the metric from the traceless sector and
determines the exchange current corresponding to that division.

\subsection{Chart data and construction of the residual tensor}
\label{vii:subsec:datos-tensor-residual}

Fix an HMT state \(x\) and a connected chart \(\mathcal U\) of
its geometric realization. On \(\mathcal U\), the nondegenerate
co-tetrad of Section~\ref{vii:sec:realizacion-cartan} determines a
smooth metric \(g=g_x\) of signature \((-+++)\). Denote its
Levi--Civita connection by \(\mathring\nabla\), and write
\[
 \mathsf G_{\mu\nu}[g]
 :=\operatorname{Ric}_{\mu\nu}[g]-\tfrac12R[g]g_{\mu\nu}.
\]
This distinguishes the Einstein tensor \(\mathsf G[g]\) from
the gravitational coupling \(G_{\rm int}\). The derivative used
in the balances of this section is \(\mathring\nabla\),
consistently with the effective metric equation of
Section~\ref{vii:subsec:fuente-metrica-efectiva}.

We work in the dimensional chart \(c_{\rm int}=1\). Assume
\(G_{\rm int}(x)>0\) and \(\Lambda_{\rm ref}(x)\) constant
on \(\mathcal U\), and set \(\kappa_x=8\pi G_{\rm int}(x)\).
The symbol \(x\) identifies the state supplying these readers;
covariant differentiation acts on the points of \(\mathcal U\).
The genealogy of the couplings remains in
Section~\ref{v:ant:sec:gravity}.

Let \(T_b\) be the smooth symmetric tensor expressed by the chart's
matter transducer; in the baryonic readout, it is its baryonic source.
Define
\begin{equation}
 \mathcal E^{\rm ref}_{\mu\nu}
 :=\kappa_x^{-1}\bigl(\mathsf G_{\mu\nu}[g]
                            +\Lambda_{\rm ref}g_{\mu\nu}\bigr),
 \qquad
 T^{\rm res}_{\mu\nu}:=\mathcal E^{\rm ref}_{\mu\nu}-T_{b,\mu\nu}.
 \label{vii:eq:residual-definicion}
\end{equation}
The tensor \(T^{\rm res}\) is the effective residual tensor,
also denoted \(T_{\rm osc}\) when interpreting the geometric
sector remaining to be expressed relative to \(T_b\). The
construction retains \(g_x,G_{\rm int},\Lambda_{\rm ref}\) and
\(T_b\) as its effective arguments.

\begin{theorem}[Uniqueness of the residual and balance with the matter source]
\label{vii:thm:residual-unicidad-balance}
For the preceding data, there is a unique symmetric tensor
\(T^{\rm res}\) satisfying
\(\mathcal E^{\rm ref}=T_b+T^{\rm res}\). Its divergence is
\begin{equation}
 \mathring\nabla^\mu T^{\rm res}_{\mu\nu}
 =-\mathring\nabla^\mu T_{b,\mu\nu}.
 \label{vii:eq:balance-residual-material}
\end{equation}
In particular, if the matter source is separately conserved,
\(\mathring\nabla^\mu T_{b,\mu\nu}=0\), then
\(T^{\rm res}\) is also conserved.
\end{theorem}
\begin{proof}
The difference of two tensors satisfying the stated equation is
zero; expression~\eqref{vii:eq:residual-definicion} supplies
existence. The contracted Bianchi identity and Levi--Civita metric
compatibility give
\[
 \mathring\nabla^\mu\mathsf G_{\mu\nu}=0,
 \qquad \mathring\nabla g=0.
\]
Since \(\kappa_x\) and \(\Lambda_{\rm ref}\) are constant
on the chart, \(\mathring\nabla^\mu\mathcal E^{\rm ref}_{\mu\nu}=0\).
Differentiating the difference defining \(T^{\rm res}\) proves
\eqref{vii:eq:balance-residual-material}, including the assertion
under separate matter conservation.
\end{proof}

Spatial and temporal constancy of the couplings on the chart fixes
the domain of this balance. A collection of states may express different
couplings in different charts; each preceding identity applies to
the chart and state specifying it.

\subsection{Canonical decomposition and exchange current}
\label{vii:subsec:traza-corriente-residual}

Define the trace function and the two tensors
\begin{equation}
 \theta:=g^{\mu\nu}T^{\rm res}_{\mu\nu},\qquad
 T^{\Lambda}_{\mu\nu}:=\tfrac14\theta g_{\mu\nu},\qquad
 T^{0}_{\mu\nu}:=T^{\rm res}_{\mu\nu}-T^{\Lambda}_{\mu\nu}.
 \label{vii:eq:descomposicion-traza-residual}
\end{equation}
The superscript \(\Lambda\) identifies the component proportional
to the metric; \(\Lambda_{\rm ref}\) is the reference parameter
of \eqref{vii:eq:residual-definicion}. Their functions are distinct.

\begin{proposition}[Trace projectors and uniqueness of the decomposition]
\label{vii:prop:proyectores-traza-residual}
On symmetric covariant tensors of order two, the maps
\[
 P_{\rm tr}(S)=\tfrac14\operatorname{tr}_g(S)g,
 \qquad P_0(S)=S-P_{\rm tr}(S)
\]
satisfy
\begin{equation}
 P_{\rm tr}^2=P_{\rm tr},\quad P_0^2=P_0,\quad
 P_{\rm tr}P_0=P_0P_{\rm tr}=0,\quad P_{\rm tr}+P_0=I.
 \label{vii:eq:proyectores-traza-residual}
\end{equation}
In particular, \(T^{\rm res}=T^\Lambda+T^0\),
\(\operatorname{tr}_gT^0=0\), and this decomposition is unique.
\end{proposition}
\begin{proof}
In four dimensions, \(g^{\mu\nu}g_{\mu\nu}=4\). Therefore,
\(\operatorname{tr}_g(P_{\rm tr}S)=\operatorname{tr}_gS\),
which proves idempotence of \(P_{\rm tr}\), vanishing of the
trace of \(P_0S\), and, by expansion, the other identities.
If \(S=fg+Q\) with \(\operatorname{tr}_gQ=0\), taking the
trace gives \(f=\operatorname{tr}_gS/4\). This also determines
\(Q=S-fg\), establishing uniqueness.
\end{proof}

\begin{theorem}[Covariant exchange between the two sectors]
\label{vii:thm:intercambio-traza-residual}
On the domain of \eqref{vii:eq:residual-definicion}, let
\(b_\nu:=\mathring\nabla^\mu T_{b,\mu\nu}\). The trace
exchange current is the one-form
\begin{equation}
 \mathcal J^{\rm tr}_\nu
 :=\tfrac14\partial_\nu\theta
 =\mathring\nabla^\mu T^\Lambda_{\mu\nu}.
 \label{vii:eq:corriente-intercambio-traza}
\end{equation}
The full balances are
\begin{equation}
 \mathring\nabla^\mu T^0_{\mu\nu}
       =-b_\nu-\mathcal J^{\rm tr}_\nu,
 \qquad
 \mathring\nabla^\mu(T^\Lambda+T^0)_{\mu\nu}=-b_\nu.
 \label{vii:eq:balances-sectores-residuales}
\end{equation}
If \(b=0\), the two sectors exchange exactly opposite currents.
In that case, they are separately conserved if and only if \(\theta\)
is constant on the connected chart.
\end{theorem}
\begin{proof}
Metric compatibility allows the calculation
\[
 \mathring\nabla^\mu\bigl(\tfrac14\theta g_{\mu\nu}\bigr)
 =\tfrac14g^{\mu\alpha}(\partial_\alpha\theta)g_{\mu\nu}
 =\tfrac14\partial_\nu\theta.
\]
Subtracting this equality from \eqref{vii:eq:balance-residual-material}
gives \eqref{vii:eq:balances-sectores-residuales}. With \(b=0\),
separate conservation is equivalent to \(\mathrm d\theta=0\);
a smooth function with zero differential is constant on every
connected component.
\end{proof}

The one-form \(\mathcal J^{\rm tr}\) measures exchange in the
tensorial division. It retains a type distinct from the transport
covector of Section~\ref{vii:sec:variacion-memoria} and the spin
current \(\sigma_{IJ}\in\Omega^3(\mathcal U;\Lambda^2E^*)\)
of \eqref{vii:eq:corriente-variacional}.

\subsection{Link to the effective Cartan--Holst source}
\label{vii:subsec:residual-fuente-cartan}

On a solution of \eqref{vii:eq:ecuacion-metrica-efectiva}, in the
same units \(c=1\), the residual has expression
\begin{equation}
 T^{\rm res}
 =T^{\rm phys,LC}+U^{\rm phys,spin}-T_b
   +\frac{\Lambda_{\rm ref}-\Lambda_{\rm int}}{\kappa_x}\,g.
 \label{vii:eq:residual-composicion-cartan}
\end{equation}
Indeed, substituting
\(\mathsf G+\Lambda_{\rm int}g
  =\kappa_x(T^{\rm phys,LC}+U^{\rm phys,spin})\)
into \eqref{vii:eq:residual-definicion} gives equality term by term.
When the matter transducer expresses precisely
\(T_b=T^{\rm phys,LC}\) and \(\Lambda_{\rm ref}=\Lambda_{\rm int}\),
one obtains \(T^{\rm res}=U^{\rm phys,spin}\).
The condition of separate conservation of \(T_b\) in
Theorem~\ref{vii:thm:residual-unicidad-balance} retains its own
function: \eqref{vii:eq:balance-metrico-total} establishes balance
of the total source, and the division among its components depends
on the matter action being used.

\subsection{Trace sector and acceleration of the selected charts}
\label{vii:subsec:residual-aceleracion-ds}

The TRIT reader of state \(x\) expresses \(s=\tau(x)\in\{+1,0,-1\}\),
and its positive clock expresses
\begin{equation}
 H_*=\frac{2\pi}{108t_0(x)}>0.
 \label{vii:eq:reloj-carta-ds}
\end{equation}
Let \(h_s\) be a three-dimensional metric of constant sectional
curvature \(s\). The three charts of the selected realization
are written
\begin{equation}
 g_s=-\mathrm dt^2+a_s(t)^2h_s,\qquad
 a_+(t)=H_*^{-1}\cosh(H_*t),\quad
 a_0(t)=\exp(H_*t),\quad
 a_-(t)=H_*^{-1}\sinh(H_*t).
 \label{vii:eq:cartas-ds-residual}
\end{equation}
The closed and flat charts are considered for \(t\in\mathbb R\);
the open chart for \(t>0\). The sign \(s\) and the clock \(H_*\)
retain their antecedents in the state selecting the realization.

\begin{lemma}[Curvature of the three charts of the same clock]
\label{vii:lem:curvatura-cartas-ds}
For the functions in \eqref{vii:eq:cartas-ds-residual},
\begin{equation}
 \frac{\ddot a_s}{a_s}=H_*^2,\qquad
 \left(\frac{\dot a_s}{a_s}\right)^2+\frac{s}{a_s^2}=H_*^2,
 \qquad
 \operatorname{Ric}[g_s]=3H_*^2g_s.
 \label{vii:eq:curvatura-cartas-ds}
\end{equation}
Consequently, \(R[g_s]=12H_*^2\) and
\(\mathsf G[g_s]=-3H_*^2g_s\).
\end{lemma}
\begin{proof}
The first two equalities follow by differentiating the three functions:
the closed and open charts use
\(\cosh^2z-\sinh^2z=1\), and the flat chart uses
\(\dot a_0=H_*a_0\). To verify the tensorial identification,
the connection coefficients involving the temporal coordinate are
\[
 \Gamma^0{}_{ij}=a_s\dot a_s(h_s)_{ij},\qquad
 \Gamma^i{}_{0j}=(\dot a_s/a_s)\delta^i_j,
\]
and the purely spatial ones agree with those of \(h_s\).
Contracting the curvature gives
\[
 \operatorname{Ric}_{00}=-3\ddot a_s/a_s,\qquad
 \operatorname{Ric}_{0i}=0,\qquad
 \operatorname{Ric}_{ij}
  =(a_s\ddot a_s+2\dot a_s^2+2s)(h_s)_{ij}.
\]
Substituting the first two identities of
\eqref{vii:eq:curvatura-cartas-ds} gives
\(\operatorname{Ric}_{00}=-3H_*^2\) and
\(\operatorname{Ric}_{ij}=3H_*^2a_s^2(h_s)_{ij}\).
Taking the trace and forming the Einstein tensor proves the other
assertions, with the curvature convention of
\eqref{vii:eq:ecuacion-metrica-efectiva}.
\end{proof}

\begin{theorem}[Density, pressure, and sign of the selected acceleration]
\label{vii:thm:residual-ds-aceleracion}
In the preceding charts, for \(\Lambda_{\rm ref}=0\) and \(T_b=0\),
the residual tensor is
\begin{equation}
 T^{\rm res}=-\rho_*g_s,\qquad
 \rho_*:=\frac{3H_*^2}{8\pi G_{\rm int}}>0,
 \qquad T^\Lambda=T^{\rm res},\quad T^0=0.
 \label{vii:eq:residual-ds-puro}
\end{equation}
For a unit comoving observer \(u=\partial_t\), the density and
isotropic pressure are \(\rho_*\) and \(p_*=-\rho_*\).
The acceleration equation evaluates to
\begin{equation}
 \frac{\ddot a_s}{a_s}
 =-\frac{4\pi G_{\rm int}}3(\rho_*+3p_*)=H_*^2>0.
 \label{vii:eq:aceleracion-residual-ds}
\end{equation}
\end{theorem}
\begin{proof}
The preceding lemma and \eqref{vii:eq:residual-definicion} give
\(T^{\rm res}=-(3H_*^2/\kappa_x)g_s\).
Since \(g_s(u,u)=-1\),
\(T^{\rm res}_{\mu\nu}u^\mu u^\nu=\rho_*\).
The spatial projector \(q^{\mu\nu}=g_s^{\mu\nu}+u^\mu u^\nu\)
satisfies \(q^{\mu\nu}(g_s)_{\mu\nu}=3\), and hence
\[
 p_*:=\tfrac13q^{\mu\nu}T^{\rm res}_{\mu\nu}=-\rho_*.
\]
The trace is \(-4\rho_*\), so
\eqref{vii:eq:descomposicion-traza-residual} produces the two
sectors in \eqref{vii:eq:residual-ds-puro}. Finally,
\(\rho_*+3p_*=-2\rho_*\), and substitution into the right-hand
side of \eqref{vii:eq:aceleracion-residual-ds} gives
\(\kappa_x\rho_*/3=H_*^2\), equal to the geometric value
already calculated. Positivity uses exactly \(G_{\rm int}>0\)
and \(H_*>0\).
\end{proof}

In this realization, the trace sector expresses a constant density
and an opposite pressure, and the internal clock fixes the positive
acceleration. In a general chart, \(T^\Lambda\) describes the
pointwise isotropic component, while \(T^0\) retains the
traceless residual, and both obey the exchange in
\eqref{vii:eq:balances-sectores-residuales}. Identifying these
sectors with observational dark sources requires a single realization
to reproduce expansion, lensing, rotation, structure growth, and the
cosmic background jointly. The present construction fixes the tensor,
its balances, and the acceleration of the selected charts entering
that comparison.
